% !TEX root = ../auv-thesis.tex
% Authoritative thesis source; edit this TeX file directly.
% Legacy Markdown material is reference-only and must not overwrite this file.

\chapter{不确定性感知的任务决策与安全控制}

\section{引言}\label{sec:ch04-4-1}

第 2 章给出了系统的硬件骨架，第 3 章解决了``AUV 在哪里、对自己的位置有多大把握''的问题。本章接着回答下一个问题：在已有状态估计的基础上，如何把``去哪里巡检''与``怎样安全到达''这两件事拆解成可工程化的决策与控制结构？在展开之前，必须先澄清两个被工程实践反复证明的基本认识。其一，AUV 在水平方向上是欠驱动的，它没有横向执行器，因而控制器无法直接命令横向位移，只能通过航向调节间接实现横向靠拢；其二，水下巡检任务实际上由任务调度、轨迹生成、参考生成、低层执行与安全仲裁五类语义完全不同的负载构成，它们的时间尺度、确定性程度和失败代价各不相同，把它们塞进同一个控制器既不现实也不安全。本章的全部设计都建立在这两个认识之上。

\subsection{欠驱动 AUV 运动约束分析}\label{sec:ch04-4-1-1}

本文考虑的 AUV 采用鱼雷型布局：尾部一个主推进器配合十字形舵面，水平舵与垂直舵分别调节俯仰和航向。这种构型在水平面内能够通过舵角控制航向和前进速度，但在横移方向上没有直接执行器，属于典型的欠驱动系统；垂直面内则通过水平舵控制俯仰角、间接改变深度。这一约束带来一个根本性的推论：控制器不能指望在横移方向上产生直接作用力，凡是``沿电缆中心线保持横向偏移为零''的目标，最终都必须被翻译成``调整航向使 AUV 朝中心线靠拢''。这个翻译过程，正是后文视线（LOS）制导与制导层模型预测控制存在的根本动机\cite{breivik2005guidance,fossen2015los}。

在仿真环境中，AUV 接受一个多维速度指令，经内层映射为垂直舵、水平舵与推进器转速；执行器本身具有一阶惯性，舵角与推力的实际响应滞后于指令，时间常数在百毫秒量级。这一滞后看似细节，却直接约束了预测控制的时域下限：如果预测时域过短，模型还来不及``感受到''舵效就要给出最优解，预测的意义便被掏空。据此，本文把模型预测控制的预测时域取为 4 秒，以充分覆盖一阶惯性的过渡过程。

控制对象的状态向量精简保留六个变量：

\begin{equation}
\mathbf{X} = [x,\ y,\ z,\ \psi,\ u,\ w]^{\mathsf{T}}
\label{eq:ch04-state-vector}
\end{equation}

其中 \((x,y)\) 为水平面位置，\(z\) 为深度（正向下），\(\psi\) 为航向角，\(u\) 为前向速度（surge），\(w\) 为垂向速度（heave）。横向速度（sway）不作为受控量，而是通过航向偏转间接影响水平轨迹，这正是欠驱动约束在状态向量层面的直接体现。俯仰角不作为独立状态量进入预测模型，而是按目标深度误差近似生成；这一取舍牺牲了俯仰动力学的精确性，却显著降低了优化问题的维度，使求解时间维持在控制周期可接受的范围内。

\subsection{决策与控制系统的分层解耦逻辑}\label{sec:ch04-4-1-2}

理解了欠驱动约束之后，下一个问题是如何在系统层面组织决策与控制的不同职责。海缆巡检任务需要在同一个系统里处理``去哪里巡检''\,``如何沿电缆走''\,``感知不可靠时如何降级''和``出现安全风险时如何退出''四类问题，而这四类问题的时间尺度、确定性和失败代价完全不同，硬塞进单一控制器只会让系统不可维护、故障不可定位。本文据此采用分层解耦设计，把系统自上而下划分为五层：任务调度层决定当前应处于搜索、锁定、跟踪、上浮、返航还是急停哪个阶段；轨迹生成层根据任务阶段产生参考路径与参考速度；参考生成层把参考路径转化为航向、深度、速度指令；低层控制层把参考指令落实为舵角与推力；安全仲裁层则始终监控感知、通信与执行器状态，必要时限制或回退输出。

\begin{longtable}[]{@{}
  >{\raggedright\arraybackslash}p{(\linewidth - 4\tabcolsep) * \real{0.3333}}
  >{\raggedright\arraybackslash}p{(\linewidth - 4\tabcolsep) * \real{0.3333}}
  >{\raggedright\arraybackslash}p{(\linewidth - 4\tabcolsep) * \real{0.3333}}@{}}
\caption{决策与控制系统分层职责}\label{tab:ch04-control-layers}\tabularnewline
\toprule\noalign{}
层级 & 职责 & 实现思路 \\
\midrule\noalign{}
\endhead
\bottomrule\noalign{}
\endlastfoot
任务调度层 & 决定当前处于搜索/锁定/跟踪/上浮/返航/急停哪个阶段 & 行为树（Behavior Tree） \\
轨迹生成层 & 根据任务阶段产生参考路径与参考速度 & 之字形扫描 + 电缆跟踪轨迹 \\
参考生成层 & 把参考路径转化为航向/深度/速度指令 & 视线制导 / 制导层模型预测控制 \\
低层控制层 & 把参考指令转化为舵角与推力 & 内层深度-航向自动驾驶（PID） \\
安全仲裁层 & 监控感知、通信与执行器状态，必要时限制或回退输出 & 指令守卫 + 失效保护节点 \\
\end{longtable}

这种分层在工程上带来两个直接好处：一是允许各层独立验证，任务调度可以在不连接控制器的情况下用模拟动作验证调度逻辑，控制器也可以脱离行为树单独做控制扫描；二是允许各层独立替换，用 PID 替换模型预测控制不影响行为树，用行为树替换状态机也不影响控制层。任务运行阶段及其准入、回退关系如图~\ref{fig:ch04-mission-state-machine}所示。该状态机描述的是任务生命周期的阶段管理；每个阶段内部的行动选择仍由后文的行为树承担，二者在职责上互补而非替代。

\begin{figure}[tbp]
\centering
\includegraphics[width=\linewidth,height=0.44\textheight,keepaspectratio,alt={IDLE、PREFLIGHT、SHADOW、SINGLE LOOP和FULL AUTONOMY的分级释放、任务完成与安全回退，任一阶段保留回退通道}]{architecture/auv_v2_mission_state_machine.pdf}
\caption{巡检任务生命周期状态机：IDLE \(\rightarrow\) PREFLIGHT \(\rightarrow\) SHADOW \(\rightarrow\) SINGLE LOOP \(\rightarrow\) FULL AUTONOMY 的分级释放、任务完成与安全回退；任一阶段保留回退通道}\label{fig:ch04-mission-state-machine}
\end{figure}

\section{基于行为树的任务规划与调度}\label{sec:ch04-4-2}

确定分层架构之后，第一个需要具体落实的是任务调度层。工程实践中可选的实现方式主要有两类，传统的有限状态机与现代的行为树。有限状态机直接、易写、易理解，但对``应急优先级抢占''和``任务可组合''的支持极为薄弱：每新增一种故障类型，都要在所有相关状态之间补画新的转移边，状态数很快膨胀到平方量级。行为树则把``组合''作为第一等公民，应急行为可以以高优先级子树的形式插入到任意位置而不破坏主任务逻辑\cite{colledanchise2018behavior,colledanchise2017how,sprague2018improving}。本文据此选择行为树作为任务调度层的实现，其``安全优先、任务后置''的结构如图~\ref{fig:ch04-behavior-tree-core}所示。

\begin{figure}[tbp]
\centering
\includegraphics[width=\linewidth,height=0.52\textheight,keepaspectratio,alt={根选择器依次检查紧急上浮、待命授权和调试级联，主任务完成下潜后由海底安全限制器约束置信度路由}]{architecture/auv_v2_behavior_tree_illustration.pdf}
\caption{行为树决策核心：根选择器按紧急、待命、任务三级优先级逐 tick 检查；主任务经下潜和海底安全限制后，在精准巡检与低置信度搜索之间路由}\label{fig:ch04-behavior-tree-core}
\end{figure}

\subsection{行为树节点设计}\label{sec:ch04-4-2-1}

行为树由控制节点、条件节点、行动节点和装饰器组成。实际实现中的根节点是 \texttt{memory=False} 的选择器，每个 tick 均从左到右重新检查三个分支。第一分支为紧急顺序节点：漏水、低电或穿底任一条件成立时，\texttt{EmergencyCondition} 返回成功并进入 \texttt{EmergencySurface}。第二分支 \texttt{StandbyCheck} 处理未获自主授权或执行脑链路不可用的情形，写入 IDLE 安全目标并返回成功，从而阻止后续任务分支运行。只有前两级均不接管时，第三分支才按调试等级选择保持、解析轨迹或主任务。

主任务不是``搜索--锁定--跟踪''三个顺序动作的硬编码展开，而是先执行 \texttt{DiveToDepth}，再由 \texttt{SeabedSafetyLimiter} 包裹路线选择器。定位辅助可用且置信度越过迟滞门限时，\texttt{PreciseInspectionSequence} 进入并行跟踪；异常状态由 \texttt{AnomalySpeedLimiter} 把速度缩放为 0.4。条件失败时路线选择器回落到 \texttt{ZigZagSearch}。因此，图~\ref{fig:ch04-behavior-tree-core}表达的是代码中可核对的优先级与路由关系，任务阶段的长期演化仍由上一节生命周期状态机描述。关键节点如下：

\begin{longtable}[]{@{}
  >{\raggedright\arraybackslash}p{(\linewidth - 4\tabcolsep) * \real{0.3333}}
  >{\raggedright\arraybackslash}p{(\linewidth - 4\tabcolsep) * \real{0.3333}}
  >{\raggedright\arraybackslash}p{(\linewidth - 4\tabcolsep) * \real{0.3333}}@{}}
\caption{行为树关键条件与行动节点}\label{tab:ch04-bt-nodes}\tabularnewline
\toprule\noalign{}
节点类型 & 功能 & 职责 \\
\midrule\noalign{}
\endfirsthead
\toprule\noalign{}
节点类型 & 功能 & 职责 \\
\midrule\noalign{}
\endhead
\bottomrule\noalign{}
\endlastfoot
条件 & \texttt{EmergencyCondition} & 漏水、低电或穿底任一成立时触发紧急分支 \\
行动 & \texttt{EmergencySurface} & 写入高优先级紧急上浮目标 \\
行动 & \texttt{StandbyCheck} & 未授权或执行脑链路不可用时写入 IDLE 目标 \\
条件 & \texttt{DebugLevel}\newline\texttt{Condition} & 将调试等级 1、2 与主任务等级 0/3 分流 \\
行动 & \texttt{DiveToDepth} & 在路线任务前建立 4 m 目标深度 \\
条件 & \texttt{ConfidenceAbove}\newline\texttt{Threshold} & 检查定位辅助可用性并通过迟滞置信度门 \\
装饰器 & \texttt{SeabedSafety}\newline\texttt{Limiter} & 对路线输出施加穿底保护和 0.5 降速 \\
装饰器 & \texttt{AnomalySpeed}\newline\texttt{Limiter} & 对异常状态下的并行跟踪输出施加 0.4 降速 \\
行动 & \texttt{ZigZagSearch} & 精准巡检条件失败时执行回落搜索 \\
\end{longtable}

\subsection{巡检任务逻辑流}\label{sec:ch04-4-2-2}

上一节给出了生产行为树的节点和优先级。本节进一步说明运行中的任务逻辑，但不把概念上的搜索、建立观测和稳定跟踪阶段误写成代码中不存在的四状态行为树。AUV 入水后可能偏离电缆先验，需要先以之字形扫描扩大覆盖；观测刚出现时噪声和不确定度仍高，也不能立即把检测结果当作连续跟踪许可。生产实现因此把这类阶段演化压缩为每个 tick 都重新判定的置信度路由：定位辅助可用且置信度通过迟滞门时进入精准巡检，否则执行之字形搜索。

精准巡检分支并行更新电缆路由、速度与离底目标；异常状态由速度装饰器将目标速度缩放为 0.4，穿底风险则由海底安全装饰器进一步限速。漏水、低电或穿底条件触发根节点的紧急上浮分支，授权撤销或执行脑链路不可用则触发待命分支。由此，搜索与跟踪的切换由观测质量驱动，应急与待命能够在下一次 tick 抢占路线任务；图~\ref{fig:ch04-mission-state-machine}中的 PREFLIGHT、SHADOW、SINGLE LOOP 和 FULL AUTONOMY 用于表达系统能力的分级释放，不等同于行为树内部节点。

这里没有单独硬编码``锁定''状态。观测从首次出现到稳定可用的过渡由第 3 章的滤波、分源质量评价和迟滞置信度门共同实现；任务层只消费其输出。这样的职责划分使行为树保持短小，并避免同一置信度状态在估计器、状态机和行为树中重复维护。

\subsection{应急自救行为的优先级机制}\label{sec:ch04-4-2-3}

正常任务流程之上，系统需要在不同执行层分配故障处置权。ROS2 行为树把漏水、低电和穿底合并为同一个最高优先级或条件；三者在该节点内没有再次排序，任一成立都会进入紧急上浮。自主授权撤销或执行脑链路不可用由第二优先级的待命节点生成 IDLE 目标。短时感知退化则不进入紧急子树，而是通过置信度路由、搜索回落和速度限制逐级处理。更底层的命令新鲜度、执行器回零和看门狗属于仲裁器与 PC104/VxWorks 的职责，不能画成行为树内部并不存在的兄弟节点。

\begin{longtable}[]{@{}
  >{\raggedright\arraybackslash}p{(\linewidth - 4\tabcolsep) * \real{0.3333}}
  >{\raggedright\arraybackslash}p{(\linewidth - 4\tabcolsep) * \real{0.3333}}
  >{\raggedright\arraybackslash}p{(\linewidth - 4\tabcolsep) * \real{0.3333}}@{}}
\caption{应急故障触发条件与响应行为}\label{tab:ch04-emergency-priority}\tabularnewline
\toprule\noalign{}
故障类型 & 触发条件 & 响应行为 \\
\midrule\noalign{}
\endfirsthead
\toprule\noalign{}
故障类型 & 触发条件 & 响应行为 \\
\midrule\noalign{}
\endhead
\bottomrule\noalign{}
\endlastfoot
行为树紧急分支 & 漏水、低电或穿底 & 生成高优先级紧急上浮目标 \\
行为树待命分支 & 自主状态非 ACTIVE 或执行脑链路不可用 & 生成 IDLE 安全目标并阻止任务分支运行 \\
任务路线分支 & 定位辅助不可用或置信度未过门限 & 从精准巡检回落到之字形搜索 \\
控制装饰器 & 穿底风险或异常状态 & 对路线/跟踪目标限速，必要时由上层紧急分支接管 \\
ROS2 仲裁与执行脑 & 指令过期、授权撤销或本地看门狗触发 & 停止释放自主指令、回零或切换本地保护 \\
\end{longtable}

这些路径形成``任务降级--自主仲裁--执行脑保护''的分层安全链。行为树负责根据已汇总的传感状态选择目标行为；ROS2 仲裁层检查授权、消息新鲜度和故障字；PC104/VxWorks 保留不依赖 Jetson 调度周期的本地看门狗与防触底保护。各守卫的输入、处置归属和目标安全状态见第 2 章图~\ref{fig:ch02-emergency-arbitration}。行为树与有限状态机的对照实验使用相同触发和去抖规则，以免把阈值整定差异误当成架构差异；第~\ref{sec:ch04-bt-fsm-figures}节给出同输入下的结果。

\subsubsection{行为树 vs 状态机对比图组}\label{sec:ch04-bt-fsm-figures}

下列实验分别从反应时延、抖动行为、复杂度规模、扩展开销和组合故障生存率五个维度对比两种方案。漏水故障在两种架构中均保持不滤波直通；1000 次持续故障注入得到相同的 100 ms 响应延迟和零遗漏，因而反应时延不构成架构差异。对置信度切换则采用完全相同的中心阈值 0.7、迟滞带宽 0.08 和 3 帧连续确认，避免用不同规则制造架构优势。在 100 条、每条 30 s 的阈值噪声序列中，单阈值方案平均发生 149.09 次模式切换，施密特触发与去抖后降至 2.23 次，降幅约 98.5\%；500 次低电压、漏水和近底穿越组合故障中，两种架构均无遗漏或死锁，任务生存率一致。因此，实验不支持``行为树响应更快''，其优势应限定为高优先级安全逻辑的集中表达和新增故障时较低的维护开销，下面仅保留能体现该优势的复杂度与扩展开销两图；反应时延、生存率与去抖降幅（98.5\%）的等价性已在上文以数值给出，不再单独作图。

\begin{figure}[htbp]
\centering
\includegraphics[width=0.95\linewidth,height=\textheight,keepaspectratio,alt={行为树与有限状态机的结构复杂度对比（圈复杂度与节点规模）}]{experiments/decision_bt_vs_fsm/03_complexity_comparison.pdf}
\caption{行为树与有限状态机的结构复杂度对比（圈复杂度与节点规模）}\label{fig:ch04-bt-fsm-complexity}
\end{figure}
\FloatBarrier

\begin{figure}[htbp]
\centering
\includegraphics[width=0.95\linewidth,height=\textheight,keepaspectratio,alt={新增故障分支的扩展开销对比（有限状态机需补转移边，行为树仅需插入高优先级子树）}]{experiments/decision_bt_vs_fsm/04_expansion_cost.pdf}
\caption{新增故障分支的扩展开销对比（有限状态机需补转移边，行为树仅需插入高优先级子树）}\label{fig:ch04-bt-fsm-expansion}
\end{figure}
\FloatBarrier

\section{不确定感知下的轨迹生成策略}\label{sec:ch04-4-3}

任务调度层决定了``当前应该做什么''，下一个问题是``按什么样的几何路径去做''。这一问题可以进一步拆成两个互补的子问题：其一，搜索阶段需要一条能够覆盖未知区域的扫描轨迹；其二，跟踪阶段在感知质量发生变化时，需要一种能在``激进巡检''与``保守等待''之间动态切换的航速规划机制。前者关注空间几何，后者关注时间动态，两者共同决定 AUV 在不同感知条件下走出什么样的物理路径。

\subsection{对齐行业标准的之字形空间扫描轨迹设计}\label{sec:ch04-4-3-1}

之字形扫描是海底巡检常用的预扫描模式\cite{choset1998boustrophedon,yordanova2020coverage}，其关键设计量不是单独的声呐波束角，而是声、磁两类传感器均能维持有效观测的联合扫描宽度。设 AUV 离底高度为 \(h\)，声呐横向波束角为 \(\beta\)，则理想几何声呐带宽为 \(W_{\mathrm{s}}=2h\tan(\beta/2)\)；将磁传感器在给定检测判据下的有效横向宽度记为 \(W_{\mathrm{m}}\)，联合扫描采用
\begin{equation}
W_{\mathrm{eff}}=\min(W_{\mathrm{s}},W_{\mathrm{m}}),\qquad
d=\lambda W_{\mathrm{eff}},\qquad
\eta=1-\lambda ,
\label{eq:ch04-coverage-spacing}
\end{equation}
其中 \(d\) 为相邻航线间距，\(\lambda\) 为间距比，\(\eta\) 为名义重叠率。取较小宽度可以防止某一通道的较大视场掩盖另一通道的覆盖缺口。

这里需要区分两类之字形：本节讨论 SEARCH 阶段的大范围覆盖扫描；TRACK / REACQUIRE 阶段则采用局部受控横切来补充磁几何可观测性。后者利用峰值时刻定位横向中心，并以半高全宽反演垂直距离，其几何作用见图~\ref{fig:ch02-active-crossing}，本节不重复插图。

为把航迹偏差纳入间距设计，设横向跟踪误差界为 \(e=\epsilon W_{\mathrm{eff}}\)。相邻航线向不利方向偏移时，每条扫描带可保证的宽度缩减为 \(W_{\mathrm{g}}=\max(W_{\mathrm{eff}}-2e,0)\)，由此得到归一化保证覆盖率
\begin{equation}
C_{\mathrm{g}}
=\min\left(1,\frac{W_{\mathrm{g}}}{d}\right)
=\min\left(1,\frac{\max(1-2\epsilon,0)}{\lambda}\right).
\label{eq:ch04-guaranteed-coverage}
\end{equation}
式\eqref{eq:ch04-guaranteed-coverage}把``两成重叠裕度''转化为可检验条件：当横向误差界为有效宽度的 \(10\%\) 时，\(\lambda=0.8\) 恰好仍满足 \(C_{\mathrm{g}}=1\)。

覆盖安全性还需与航程和时间共同评价。对沿航线方向长度为 \(L\)、横跨方向宽度为 \(B\) 的矩形区域，航线数取
\begin{equation}
n=\left\lceil\frac{B-W_{\mathrm{eff}}}{d}\right\rceil+1 ,
\label{eq:ch04-line-count}
\end{equation}
并以直线扫描段加直径为 \(d\) 的半圆换向近似总航程 \(S=nL+(n-1)\pi d/2\)，任务时间为 \(T=S/v\)。本文以 \(W_{\mathrm{eff}}\) 的参考尺度归一化区域、航程和速度，对声呐波束角、磁有效宽度、间距比、速度比及横向误差比执行了 \(4\times4\times5\times4\times4=1280\) 组全因子解析扫描。表~\ref{tab:ch04-normalized-coverage}给出参考波束、等宽声磁通道和 \(\epsilon=0.10\) 下的间距对照。

\begin{longtable}[]{@{}rrrrrr@{}}
\caption{横向误差比为 0.10 时的归一化覆盖与时间权衡}\label{tab:ch04-normalized-coverage}\tabularnewline
\toprule\noalign{}
间距比 & 重叠率 & 保证覆盖率 & 航线数 & 归一化航程 & 时间指数 \\
\midrule\noalign{}
\endfirsthead
\toprule\noalign{}
间距比 & 重叠率 & 保证覆盖率 & 航线数 & 归一化航程 & 时间指数 \\
\midrule\noalign{}
\endhead
\bottomrule\noalign{}
\endlastfoot
0.6 & 40\% & 100.0\% & 33 & 1350.16 & 1.311 \\
0.7 & 30\% & 100.0\% & 29 & 1190.79 & 1.156 \\
0.8 & 20\% & 100.0\% & 25 & 1030.16 & 1.000 \\
0.9 & 10\% & 88.9\% & 23 & 951.10 & 0.923 \\
1.0 & 0\% & 80.0\% & 20 & 829.85 & 0.806 \\
\end{longtable}

\begin{figure}[htbp]
\centering
\includegraphics[width=0.88\linewidth,height=\textheight,keepaspectratio,alt={横向误差比为百分之十时扫描间距、保证覆盖率与任务时间的归一化权衡}]{experiments/coverage/normalized_coverage_time_tradeoff.pdf}
\caption{扫描间距对保证覆盖率与归一化任务时间的影响}\label{fig:ch04-normalized-coverage-tradeoff}
\end{figure}
\FloatBarrier

结果表明，\(\lambda=0.8\) 是当前离散设计集中仍保持完全保证覆盖的最大间距；增至 0.9 后保证覆盖率降为 88.9\%，而减至 0.7 和 0.6 虽继续保持完全覆盖，却分别增加约 15.6\% 和 31.1\% 的任务时间。速度只按反比改变任务时间而不改变几何覆盖，波束角与磁有效宽度则通过二者的较小值共同决定航线数量。因此，0.8 倍联合有效宽度不是孤立经验常数，而是在 10\% 归一化横向误差界下的覆盖--时间折中点。这里给出的均为无量纲设计规律；实际米制间距与任务时长仍须由实测波束图、作业高度、磁检测宽度、最小转弯半径和速度上限共同标定。

\subsection{基于感知置信度的动态航速规划模型}\label{sec:ch04-4-3-2}

空间扫描的几何确定之后，下一个问题是时间维度上的速度调节。在巡检阶段，AUV 不应当总以恒定速度前进，感知置信度变化时，行驶节奏也应随之变化。一个朴素的直觉是：感知质量稳定时可以按设计速度沿电缆巡检，感知置信度下降时则应降低航速、增加安全裕度。把这一直觉量化的核心，是将感知不确定性映射为速度上限的连续调节函数，而不是仅在``正常/异常''两种离散状态之间硬切换。

为什么不用硬切换而要用连续映射？原因有二。其一，硬切换会在阈值附近产生抖动，当置信度在阈值上下小幅波动时，速度上限会反复跳变，导致 AUV 不停加减速；其二，连续映射允许``轻度退化时小幅减速、严重退化时大幅减速''，把可用的感知信息按比例利用，而不是把所有不可靠观测粗暴地一律降级。这一机制与后文的不确定性感知控制权重调整深度耦合，其运作可归纳为三点。

首先是感知置信度到速度上限的映射。系统维护一个由状态协方差换算得到的实时置信度；质量控制策略可随置信度下降压低参考速度，使运动节奏与可用观测相匹配。其次是与控制权重联动：跟踪权重在低置信度区间放大，而控制努力惩罚经 sigmoid 平滑降低，为必要纠偏保留控制权威；指令平滑由速率和带宽约束保证。最后是任务路由：生产行为树采用可配置的中心阈值，定位辅助不可用或门控不通过时由精准巡检回落到之字形搜索。紧急上浮只由漏水、低电和穿底等独立安全条件触发，不由低置信度单独触发。协方差到置信度的映射见第 3 章，实际授权和速度缩放证据见第~\ref{sec:ch05-5-5-6}节。

\section{不确定性感知模型预测控制设计}\label{sec:ch04-4-4}

任务调度层与轨迹生成层就位之后，本章进入最关键的一层，也就是本文最主要的算法贡献：不确定性感知的模型预测控制（Uncertainty-Aware MPC，下称 UA-MPC）。在具体构造上，其约束设置与制导层定位受到约束型 MPC 和 AUV 轨迹跟踪研究的启发\cite{mayne2000constrained,li2021trajectory}；区别在于，本文进一步把在线感知置信度引入代价权重。进入具体公式之前，需要先澄清一个工程上的定位问题：本系统中的 MPC 并不是直接驱动舵机的低层控制器，而是\textbf{制导层的参考生成器}。作出这一定位有两条理由。其一，内层已经有一个调优良好的深度-航向 PID 自动驾驶，强行用 MPC 替代它不仅得不到额外收益，反而引入不必要的求解时间；其二，MPC 的真正优势在于多变量耦合、约束处理、未来预瞄和不确定性权重，只有在制导层才能充分体现，而非在单一深度闭环里。本节按``模型、代价、不确定性权重、参考、约束、求解器''六个层次依次展开。

\subsection{离散化运动学预测模型}\label{sec:ch04-4-4-1}

按上述路线，本节首先回答``模型层''应当采用什么样的状态空间表示。MPC 是一种模型驱动的控制方法，预测模型是它的全部物理依据，模型设计是否合理直接决定后续代价函数与约束的讨论是否有意义。这里从两个问题入手：状态向量应取多少维？离散化应采用哪种数值方案？

先看状态维度的选择。全自由度刚体模型把运动分解为三个平动与三个转动，连同位形和速度共十二维。但本系统的 MPC 不需要承担这一全套，理由有三：其一，欠驱动 AUV 在水平面没有横向执行器，横向速度由航向偏转间接产生，把它作为独立状态量不会带来控制自由度，反而使求解维度翻倍；其二，横滚在巡检任务中本就被内层的小角度控制抑制，纳入预测模型只会让动力学方程复杂化；其三，俯仰角由目标深度与当前深度之差经限幅近似生成，作为派生量进入垂向通道即可，无需独立状态。综合这三点，本系统把 MPC 状态向量精简为六维：

\begin{equation}
\mathbf{X} = [x,\ y,\ z,\ \psi,\ u,\ w]^{\mathsf{T}}
\label{eq:ch04-mpc-state}
\end{equation}

其中 \((x,y)\) 为水平面位置，\(z\) 为深度（正向下），\(\psi\) 为航向角，\(u\) 为前向速度，\(w\) 为垂向速度。这一取舍在``够用即可''与``求解时间可控''之间取得了工程平衡，既能刻画路径预瞄与深度调节的核心动力学，又把单步求解时间维持在控制周期可消化的范围。相应地，控制量精简为三维：

\begin{equation}
\mathbf{U} = [\psi_{\mathrm{cmd}},\ z_{\mathrm{cmd}},\ T_{\mathrm{cmd}}]^{\mathsf{T}}
\label{eq:ch04-mpc-control}
\end{equation}

即航向参考、深度参考与推力百分比。这再次呼应了``制导层参考生成器''的定位：MPC 输出的不是舵角和转速本身，而是三个参考量，由内层执行器在每个控制周期内消化为具体的舵面动作与转速。这种``上层只给目标、下层负责落实''的解耦既保留了调优良好的内层 PID，又把多变量耦合与约束处理交给优化器完成。

把上述状态量与控制量代入连续时间动力学，得到六个运动学方程：

\begin{equation}
\begin{aligned}
\dot{x} &= u\cos\psi, &
\dot{y} &= u\sin\psi, \\
\dot{z} &= -u\sin\theta+w\cos\theta, &
\dot{\psi} &= k_{\psi}(\psi_{\mathrm{cmd}}-\psi), \\
\dot{u} &= \frac{T_{\mathrm{actual}}-d_u u|u|}{m_u}, &
\dot{w} &= \frac{-d_w w+k_z(z_{\mathrm{cmd}}-z)+f_b}{m_w}.
\end{aligned}
\label{eq:ch04-continuous-model}
\end{equation}

其中俯仰角 \(\theta\) 由目标深度误差经比例增益限幅得到，限幅范围约为正负三十度。作这一限幅有两重工程动机：一是 AUV 实际能够稳定保持的俯仰角受控制权威与水动力恢复力矩约束，超过约三十度后舵效非线性显著加剧，预测模型赖以成立的小角度近似便会失效；二是限幅使``目标深度→俯仰角→垂向速度''这一派生链条不会因深度误差过大而产生不切实际的瞬时俯仰指令。在深度正向下的约定里，目标深度更深时对应负俯仰角；推力则经非负限幅，确保不产生反向推力，本系统主推为单向桨，物理上不允许负推力。式\eqref{eq:ch04-continuous-model}中其余系数（航向增益、阻力、质量、深度到垂速的耦合增益与浮力项等）的取值汇总于附录 A.4。

离散化采用显式欧拉向前法：

\begin{equation}
\mathbf{X}_{k+1}
= \mathbf{X}_k+\Delta t\,f(\mathbf{X}_k,\mathbf{U}_k)
\label{eq:ch04-euler-discretization}
\end{equation}

为什么用显式欧拉而非四阶龙格-库塔一类的高阶方法？答案来自实时控制的求解时间预算。整个预测时域内有二十步、每步六维状态与三维控制，合计数百个变量参与优化，求解器在每次迭代中都要构建并因式分解雅可比矩阵；离散化方案越复杂，单次构建时间越长。四阶方法把每步动力学求值次数由一次提高到四次，对雅可比的稀疏结构和符号微分时间都有数倍的放大效应，对当前的实时性要求是不可接受的。显式欧拉在 0.2 秒步长下的局部截断误差为步长平方量级，对 4 秒预测时域内的位置预测通常只有分米级偏差，已足以支撑制导层的参考生成；
真正需要精确积分的低层执行由内层 PID 以更短步长完成，无需 MPC 自身承担。据此，时间步长取 0.2 秒、预测步数取二十步、预测时域覆盖 4 秒未来状态，这一时域选择同样不是任意的，而是由``百毫秒级一阶执行器惯性需要数倍时间常数才能完整反映稳态''反推得出，与 4.1.1 节关于执行器一阶惯性的讨论一致。

\subsection{代价函数构建}\label{sec:ch04-4-4-2}

模型层确定之后，MPC 在每个控制周期求解的优化问题，其代价由跟踪代价、控制代价与终端代价三部分构成：

\begin{equation}
\min J = J_{\mathrm{tracking}}+J_{\mathrm{control}}+J_{\mathrm{terminal}}
\label{eq:ch04-cost-decomposition}
\end{equation}

跟踪代价对预测时域内每一步累加状态与参考之差的加权平方：

\begin{equation}
\begin{aligned}
J_{\mathrm{tracking}}=\sum_k \bigl[
&w_x(x_k-x_{\mathrm{ref}})^2+w_y(y_k-y_{\mathrm{ref}})^2 \\
&+w_z(z_k-z_{\mathrm{ref}})^2+w_\psi(\psi_k-\psi_{\mathrm{ref}})^2 \\
&+w_u(u_k-u_{\mathrm{ref}})^2+w_w(w_k-w_{\mathrm{ref}})^2
\bigr].
\end{aligned}
\label{eq:ch04-tracking-cost}
\end{equation}

各维权重的相对大小反映了它们对巡检任务的重要程度：近底巡检中深度跟踪权重最大，因为深度误差直接侵蚀离底安全裕度；航向次之；水平面位置与速度则相对宽松。控制代价对控制时域内每一步的指令增量施加二次惩罚，以抑制输出的剧烈波动：

\begin{equation}
J_{\mathrm{control}}
=\sum_k s_u(c)\left(
W_{\psi_{\mathrm{cmd}}}\Delta\psi_{\mathrm{cmd},k}^2
+W_{z_{\mathrm{cmd}}}\Delta z_{\mathrm{cmd},k}^2
+W_T \Delta T_{\mathrm{cmd},k}^2
\right)
\label{eq:ch04-control-cost}
\end{equation}

其中 \(\Delta u_{\cdot,k}\) 为相邻两步指令之差，惩罚指令增量而非指令幅值可直接抑制控制抖动而不额外偏置稳态工作点；整体缩放因子 \(s_u(c)\) 由不确定性感知机制根据置信度动态调整，详见 4.4.3 节。终端代价 \(J_{\mathrm{terminal}}=\mathbf{e}_N^{\mathsf T}\mathbf{Q}_f\,\mathbf{e}_N\) 对预测时域末端的状态误差 \(\mathbf{e}_N\) 施加额外加权，以改善有限时域下的收敛性。跟踪权重、控制权重与终端权重的具体默认取值汇总于附录 A.4。

\subsection{不确定性感知的权重动态调整算法}\label{sec:ch04-4-4-3}

UA-MPC 是本文的核心创新之一。在给出公式之前，先把它要解决的直觉问题讲清楚：当感知可靠时，控制器应当积极跟踪参考路径、果断给出指令；当感知变得不可靠时，控制器应当自动收敛激进程度，既在已知参考附近保守跟随、不贸然远离，又为必要的纠偏动作保留控制权威，并由速率与带宽约束抑制指令抖动。传统 MPC 的代价权重是固定的，无法表达``感知质量在变、控制风格也应随之变''这层意思。已有的鲁棒 MPC 与管道（tube）MPC 通过约束收紧来抵御有界扰动\cite{mayne2005robust,jimoh2024tube}，灵敏度感知 MPC 则以对参数不确定性的灵敏度为代价加权\cite{belvedere2025sensitivity}，但二者都不直接以在线感知置信度调制代价权重。UA-MPC 的做法，就是把感知置信度引入代价权重，使跟踪与控制两类权重都随置信度连续变化。它体现在两个相互配合的机制上。

第一个机制是跟踪权重的不确定性感知放大。当置信度下降时，水平位置与航向维度的跟踪权重被放大，使优化器更倾向于贴着已知参考走、减少偏离：

\begin{equation}
\rho(c)=(1-c)^\alpha,
\qquad
w_e(c)=w_{e,0}\left[1+s_e\rho(c)\right]
\label{eq:ch04-confidence-weight}
\end{equation}

式中 \(c\in[0,1]\) 为感知置信度，\(s_e\) 为低置信度跟踪权重的附加增益，\(\alpha\) 为置信度映射的幂指数。选择幂函数 \(\rho(c)=(1-c)^\alpha\) 而不是线性映射，是为了让权重在置信度较高时几乎不变、只有当置信度明显下降时才快速放大，从而把``保守跟随''的代价集中施加在真正不可靠的区段；当幂指数取一时，式\eqref{eq:ch04-confidence-weight}退化为线性调整。深度维度的放大系数取水平维度的一半，因为深度通道另有内层 PID 直接兜底，不需要与水平通道同等强度的保守化。

第二个机制是控制代价的平滑调整。式\eqref{eq:ch04-control-cost}中的整体缩放因子 \(s_u(c)\) 随置信度以 S 型（sigmoid）曲线过渡：

\begin{equation}
\begin{aligned}
\sigma(x)&=\frac{1}{1+\exp(-x)},\\
s_u(c)&=\gamma_u+(1-\gamma_u)\sigma\!\left(k(c-c_0)\right).
\end{aligned}
\label{eq:ch04-control-scale}
\end{equation}

式中 \(k\) 控制过渡的陡峭程度，\(c_0\) 为置信度转折点，\(\gamma_u\) 为低置信度下控制代价的最小缩放比例。置信度高于转折点时缩放因子接近一，控制指令正常输出；低于转折点时缩放因子平滑降到下限，控制努力惩罚随之减小，使优化器在跟踪权重已放大的同时仍保留必要的纠偏能力；指令平滑性则由后述速率与带宽约束保证。之所以用 S 型曲线而非硬阈值，是因为硬切换会在阈值附近让权重突变、进而引起控制量跳变，而平滑过渡消除了这一隐患。

为便于消融验证，UA-MPC 保留一个模式开关，可在``启用完整不确定性感知机制''与``忽略置信度、权重恒定的基线''之间切换，后者作为对照基线。围绕上述两个机制的可调量，放大上限、幂指数、S 型陡峭度、置信度阈值、控制缩放下限等，本文设计了一组消融变体，用于逐一剥离各机制的贡献。这些参数的默认取值与消融定义整理于附录 A.3。

实验章的消融从可解性和横向精度两个相互独立的指标验证了上述分解。在三场景、三种子的 PVS 闭环中，固定权重与关闭 S 型控制代价降权的变体均退化为全程回退；保留该降权后，默认 UA-MPC 与线性置信度映射变体的回退率显著下降。在确定性长波 S 弯中，默认 UA-MPC 使横向 RMSE 相对关闭 S 型降权时下降。两层证据共同说明，低置信度下减小控制努力惩罚、为纠偏保留控制权威，是可解性与精度改善的主要来源；把 \(\alpha\) 从 1.5 改为 1 仅在该运动学基准中带来极小附加改善，不足以支持修改默认参数。

\subsubsection{MPC 控制效果实验图组}\label{sec:ch04-mpc-response-figures}

图~\ref{fig:ch04-mpc-summary}以 2×2 汇总方式给出制导层 MPC 在深度阶跃、航向阶跃、电缆深度跟踪与电缆航向跟踪四类响应上的形态：四个通道均无超调发散、稳态收敛，验证了离散化预测模型与硬约束设计在接口层的可用性。单通道结果均为 PVS \(n=1\) 单元验证，只用于检查接口和响应形态；控制器的多场景稳健结论由第~\ref{sec:ch05-5-ctrl}节的统一实验判定，复现入口见附录。

\begin{figure}[htbp]
\centering
\includegraphics[width=0.80\linewidth,height=\textheight,keepaspectratio,alt={MPC 综合性能汇总（深度、航向与电缆跟踪四场景 2×2 汇总，PVS 仿真，\(n=1\)）}]{experiments/control_mpc/05_mpc_combined_summary.pdf}
\caption{MPC 综合性能汇总（深度、航向与电缆跟踪四场景 2×2 汇总，PVS 仿真，\(n=1\)）}\label{fig:ch04-mpc-summary}
\end{figure}
\FloatBarrier

\subsection{参考轨迹生成}\label{sec:ch04-4-4-4}

模型、代价、不确定性权重三层确定后，下一个问题是优化器``追什么''，也就是参考从哪里来。不同任务阶段的参考来源截然不同：跟踪阶段需要从已知电缆点序列中抽取一段沿路前推的预瞄路径，搜索阶段则没有路由可依、只能由当前航向与深度自行延展。把这两类语义压进同一段逻辑会让接口臃肿，本系统据此把参考生成分为两种互斥模式，由一个开关切换。

第一种是视线动态模式，用于已经获得稳定电缆路由的跟踪阶段。它基于电缆点序列与视线前视距离动态生成参考，算法沿电缆点序列逐步前推，前视距离随时间递增。这种``边走边伸前视''的设计意味着：预测时域内各步参考点对应的不是当前视线锚点的原地复制，而是沿电缆方向逐步向前的真实路径采样，从而让优化器可以同时兼顾``当前位置如何修正''和``未来若干步即将面对的弯道''。第二种是恒定航向模式，用于搜索或电缆观测尚未稳定的阶段。它不依赖电缆点先验，完全由当前航向、当前深度和参考速度推算未来若干步的状态，使 MPC 在尚未确认电缆位置时仍能工作，而不至于因``参考缺失''而求解失败。

无论哪种模式，最终输出都是一份维度与预测时域对齐的参考矩阵，每一列对应一个时间步的期望状态、与六维状态向量一一对应。把参考从``标量序列''扩展到``完整状态序列''是 MPC 的标准做法，它让代价函数可以对每个状态分量单独施加权重，而不是只对一个综合误差求平方，这一点已在 4.4.2 节的代价函数中体现。

\subsection{执行器物理饱和约束处理}\label{sec:ch04-4-4-5}

参考准备好之后，下一个问题是优化器可以``在多大范围内''自由选择控制量。MPC 在数学上可以输出任意指令值，但实际执行器存在舵角限幅、变化率限制、推力上下限等物理约束；任何超出这些约束的指令在低层都会被强制限幅，反而使 MPC 的模型预测与真实执行轨迹脱节。本节把约束体系分为``硬约束''与``速率/带宽扩展约束''两层，前者保证指令物理可执行，后者保证指令工程可平滑。

硬约束体系直接对应执行器的物理边界：

\begin{longtable}[]{@{}lll@{}}
\caption{模型预测控制硬约束体系}\label{tab:ch04-hard-constraints}\tabularnewline
\toprule\noalign{}
约束类型 & 表达式 & 物理含义 \\
\midrule\noalign{}
\endfirsthead
\toprule\noalign{}
约束类型 & 表达式 & 物理含义 \\
\midrule\noalign{}
\endhead
\bottomrule\noalign{}
\endlastfoot
航速下限 & \(u_k\ge u_{\min}\) & 确保舵效 \\
航向指令限幅 & \(\psi_{\mathrm{cmd}}\in[\psi_{\min},\psi_{\max}]\) & 航向指令合理性 \\
深度指令限幅 & \(z_{\mathrm{cmd}}\in[z_{\min},z_{\max}]\) & 深度指令合理性 \\
推力下限 & \(T_{\mathrm{cmd}}\ge T_{\min}\) & 防止推力清零失稳 \\
推力上限 & \(T_{\mathrm{cmd}}\le T_{\max}\) & 执行器饱和 \\
\end{longtable}

各项硬约束都有明确动机：航速下限来自欠驱动 AUV 的物理事实，舵面靠水流作用力产生力矩，速度太低时舵效近乎为零，控制器再多调整也无法形成有效转向力，因此必须给航速设下限；推力下限取零而非任意负值，是因为主推为单向桨、物理上不允许反推；深度上限则与本系统设计的最大下潜深度对齐。

在硬约束之外，还需要再叠加一层速率与带宽约束：

\begin{longtable}[]{@{}
  >{\raggedright\arraybackslash}p{(\linewidth - 4\tabcolsep) * \real{0.3333}}
  >{\raggedright\arraybackslash}p{(\linewidth - 4\tabcolsep) * \real{0.3333}}
  >{\raggedright\arraybackslash}p{(\linewidth - 4\tabcolsep) * \real{0.3333}}@{}}
\caption{模型预测控制速率与带宽扩展约束}\label{tab:ch04-rate-band-constraints}\tabularnewline
\toprule\noalign{}
约束类型 & 表达式 & 物理含义 \\
\midrule\noalign{}
\endfirsthead
\toprule\noalign{}
约束类型 & 表达式 & 物理含义 \\
\midrule\noalign{}
\endhead
\bottomrule\noalign{}
\endlastfoot
深度指令速率 & \(\lvert z_{\mathrm{cmd},k+1}-z_{\mathrm{cmd},k}\rvert\le\Delta z_{\max}\) & 深度指令变化率 \\
航向指令速率 & \(\lvert\psi_{\mathrm{cmd},k+1}-\psi_{\mathrm{cmd},k}\rvert\le\Delta\psi_{\max}\) & 航向指令变化率 \\
深度指令带宽 & \(\lvert z_{\mathrm{cmd},k}-z_{\mathrm{current}}\rvert\le z_b\) & 深度指令不能离当前态太远 \\
航向指令带宽 & \(\lvert\psi_{\mathrm{cmd},k}-\psi_{\mathrm{current}}\rvert\le\psi_b\) & 航向指令不能离当前态太远 \\
\end{longtable}

为什么硬约束之外还要再加这一层？答案来自实测中观察到的两类不利现象。其一，纯硬约束允许优化器在相邻两步给出相差数十度的航向指令，这种跳变会让内层 PID 频繁过冲，把``平滑参考''变成``阶跃参考''；其二，优化器可能在某一步给出远离当前态的指令，例如当前深度五米、下一步却给二十余米的目标深度，即便落在深度硬限幅之内是合法的，但这种``瞬时大步跳''无法被一阶执行器跟随。速率约束限制相邻步之间的指令斜率，带宽约束限制任一时刻指令偏离当前态的最大距离，前者管短期变化率、后者管中期偏离量，两者一起把优化器的解空间收敛到内层 PID 能稳定跟随的窗口内。这两类约束可以分别启用或禁用以便消融：早期基线仅保留硬约束体系，引入速率与带宽约束后控制指令更平滑、求解时间略有上升，这一权衡由控制平滑度对照给出量化证据。上述各约束的默认取值列于附录 A.4。

\subsection{求解器实现与热启动}\label{sec:ch04-4-4-6}

约束体系给定后，最后一步是把``如何求解这个非线性优化问题''的工程实现说清楚。MPC 在每个 50 ms 控制周期都要重新求解一次包含数百个变量的非凸非线性优化问题，求解器的选择与实现直接决定整个系统能否实时运行。本系统采用符号建模框架配合内点法求解器\cite{andersson2019casadi,waechter2006ipopt}：前者把代价函数与约束自动符号微分得到雅可比，后者以内点迭代逐次逼近最优解。要做到实时可用，工程实现上有三个关键考量。

其一是``构建一次、复用多次''。优化问题的变量数量、约束类型和雅可比稀疏模式在配置确定后固定；系统在初始化阶段完成结构编译，运行期只更新参考、当前状态和置信度。其二是热启动。每个周期把上一次最优控制序列前移一步作为初值，使求解器从相邻解附近继续搜索；闭环对照支持其作为默认配置，迭代数与耗时统计见第~\ref{sec:ch05-5-5-8}节。其三是求解失败的兜底。当求解器返回失败状态时，控制层按配置保持上一次输出或退化为开环设定点，从而维持输出连续；该路径接入 4.2.3 节的控制器异常回退规则。

回退机制只能保证输出连续，不能替代对失败原因的观测。为区分``求解困难''、``初值不连续''和``控制周期被阻塞''，本文对每个控制周期统一记录求解器返回状态与迭代数、是否实际使用热启动、相邻状态初值跳变量、平移前后控制初值的均方根变化、初始点与最终迭代点的最大约束违反量、活跃约束数量以及当前周期墙钟时间。活跃约束按距上下界不超过 \(10^{-4}\) 判定；周期墙钟时间超过控制周期时另记阻塞事件。若失败后沿用上一成功控制指令，沿用的只是物理输出，求解状态、迭代数和耗时仍取当前失败周期，不能把上一成功帧的求解时间重复写入失败帧。由此，回退率反映``有多少周期没有产生新解''，当前墙钟时间反映``失败是否阻塞实时更新''，约束残差与初值连续性则用于定位``为什么没有产生新解''，三类语义彼此分离。

在数值容差上，本系统以最大迭代次数限制单次失败对控制周期的占用。独立扫描表明，单纯放宽迭代上限会增加失败周期阻塞，而对可解性的改善有限；软约束也会增加优化维度和阻塞。因此默认配置冻结为迭代上限 100、热启动开启、软约束关闭，终端保持由任务层在路径完成后接管。完整回退率、失败耗时和软/硬约束对照集中见第~\ref{sec:ch05-5-5-8}节。此外，优化器内部、MPC 输出端和内层执行器依次施加控制量限幅，以降低单处配置误差继续传播的风险；这类冗余的作用范围仍由故障注入验证界定。求解器配置与回退语义列于附录 A.5。

\subsection{近底安全过滤与控制障碍函数}\label{sec:ch04-4-4-7}

地形跟随和电缆巡检都要求车辆在尽量靠近海床的同时保持最小安全净空。单纯把目标离底高度设为常数并不足够，因为海床估计、车辆深度响应和控制周期离散化都会引入滞后：当未来数秒内海床快速抬升时，当前时刻看似安全的深度目标可能在执行到位前变成低净空风险。为此，本文在地形跟随候选深度之后叠加一层制导级安全过滤器，用控制障碍函数（control barrier function, CBF）的思想约束深度参考，而不是直接改写内层 PID 或 MPC 求解器。

令 \(z\) 表示车辆深度（向下为正），\(S\) 表示海床深度，\(d_{\min}\) 表示最小安全净空，则当前安全裕度定义为
\begin{equation}
  h = S - z - d_{\min}.
  \label{eq:ch04-cbf-margin}
\end{equation}
当 \(h\ge 0\) 时，车辆距离海床至少保留 \(d_{\min}\) 的净空；当 \(h<0\) 时，目标深度已经越过安全边界。离散控制周期内，CBF 的核心要求不是简单把 \(h\) 保持为正，而是限制它向边界衰减的速度：
\begin{equation}
  h_{k+1} \ge (1-\alpha_{\mathrm{cbf}}) h_k,\quad 0<\alpha_{\mathrm{cbf}}\le 1 .
  \label{eq:ch04-cbf-decay}
\end{equation}
式\eqref{eq:ch04-cbf-decay}给出了一个可解释的安全缓冲：\(\alpha_{\mathrm{cbf}}\) 越大，允许越快接近边界；\(\alpha_{\mathrm{cbf}}\) 越小，参考深度越保守（此处的衰减系数 \(\alpha_{\mathrm{cbf}}\) 与式\eqref{eq:ch04-confidence-weight}中的置信度幂指数 \(\alpha\) 是两个独立参数）。

考虑地形跟随的预瞄需求，本文不只使用当前海床 \(S_k\)，还计算预测时间窗后的海床 \(S_{k+\ell}\)，并取二者中更浅、风险更高的一项作为过滤器的保守海床：
\begin{equation}
  S_{\mathrm{safe}}=\min(S_k,S_{k+\ell}) .
  \label{eq:ch04-cbf-safe-floor}
\end{equation}
候选深度 \(z_{\mathrm{ref}}\) 经过过滤后不得超过由 \(S_{\mathrm{safe}}\)、\(d_{\min}\) 和 CBF 衰减项共同给出的上界。若当前净空已经足够大，衰减项只在有限裕度内生效，避免车辆因历史大净空而长期被压在过浅深度；若当前或预瞄净空接近边界，过滤器同时限制最大下潜速率，使深度参考以车辆能够跟随的速度接近海床。

低净空状态下还需要比普通 CBF 更直接的保护。本文设置紧急净空阈值，当实际离底高度低于该阈值时，过滤器不再等待候选参考，而是强制生成上浮深度参考，并将目标航速压低；在普通接近边界但尚未进入紧急区间时，则按当前与预瞄净空裕度线性缩放 \(v_{\mathrm{ref}}\)。因此，CBF 在本文中同时作用于两个制导量：深度参考保证不继续向海床逼近，速度门控降低车辆在低净空区间内横向推进带来的响应压力。

这层设计的定位是``安全过滤''而不是``性能优化''。它不试图改善电缆横向 RMSE，也不替代跟踪授权对观测质量的判断；它只回答一个更底层的问题：在近底巡检时，候选深度是否仍位于可信安全域内。相应实验分别使用运动学参考后端和原生 PVS 深度内环，前者检验过滤逻辑的多种子稳定性，后者检验统一执行链；二者的统计和适用边界在第~\ref{sec:ch05-5-ctrl}节给出，不能合并外推为所有极端场景均已安全闭合。

\section{验证命题与实验对应关系}\label{sec:ch04-4-5}

前面各节给出了任务调度、轨迹生成、UA-MPC、约束处理与安全过滤的设计。对应实验不在本章重复报告，而由第~\ref{sec:ch05-5-ctrl}节按三个相互独立的命题组织：近底安全实验检验地形参考与 CBF 是否守住净空边界；复杂路径实验检验预测控制相对 PID 与视线制导的适用区间；不确定性实验检验置信度调制对求解可行性、控制平滑和跟踪授权的影响。三类证据使用不同场景、基线和指标，不能拼接成单一的``全面优于''结论。

\subsection{感知退化下的控制调制命题}\label{sec:ch04-4-5-1}

该命题检验低置信度是否按式~\eqref{eq:ch04-confidence-weight}和式~\eqref{eq:ch04-control-scale}改变跟踪权重与控制代价，以及这种调制是否在固定求解预算内保持可解。评价量包括 fallback、求解耗时、控制变化率和横向误差；其中控制平滑改善与横向精度改善是两个不同结论。行为树和授权逻辑仍负责在观测不可用时降级或退出，UA-MPC 不承担从失效观测中恢复状态的职责。正式分源置信矩阵、消融变体及其负结果见第~\ref{sec:ch05-5-5-6}节和第~\ref{sec:ch05-5-5-8}节。

\subsection{复杂几何与近底安全命题}\label{sec:ch04-4-5-2}

复杂路径实验比较长波/短波 S 弯、发卡和直角路径中的曲率预瞄能力，近底实验则独立比较地形参考、内层执行和安全过滤。前者以横向误差和控制代价为主，后者以离底误差、最小净空、低净空违规与穿底率为主；安全通过不能替代路径精度，离线路径优势也不能外推为原生执行链性能。公平基线、运行规模、PID 内层 profile、CBF 多种子矩阵和原生 PVS 四相重跑均集中在第~\ref{sec:ch05-5-ctrl}节，极端场景中的终端尾段与回退边界保留在同章综合讨论中。

\section{本章小结}\label{sec:ch04-4-6}

本章围绕``观测不确定时怎样决定并安全执行动作''这一问题，按``约束→任务调度→轨迹生成→控制优化→安全过滤→验证命题''组织方法。行为树与任务状态机管理优先级和生命周期；搜索覆盖与局部主动横切承担不同的观测任务；UA-MPC 把置信度映射为跟踪权重和控制代价；CBF、速率限制、仲裁器和底层保护共同约束最终执行。各层通过明确接口连接，又保留独立回退路径。

本章给出的是一套分层决策、控制与安全框架，而不是一个在所有工况下占优的控制器。近底净空、复杂路径预瞄、不确定性调制和故障回退分别对应不同验证命题；运行规模、量化结果、收益与代价集中在第~\ref{sec:ch05-5-ctrl}节，避免方法章和实验章重复呈现同一组证据。
